\section*{4 Palette-Alternating Tree Codes}
In this section we prove Theorem4. To this end we recall the definition of the (field) trace function $\operatorname{Tr}:\mathbb F_4\to\mathbb F_2$ that is given by $\operatorname{Tr}(x)=x+x^2$. Observe that the trace function is an $\mathbb F_2$-linear map whose image and kernel are $\mathbb F_2$. In particular, if $X$ is uniform over $\mathbb F_4$, then $\operatorname{Tr}(X)$ is uniform over $\mathbb F_2$.
Let $\epsilon$ be a given parameter and define $b=\lceil1/\epsilon\rceil$. Let $\{R_i\}_{i\in\mathbb N}$ be a sequence of independent random variables, each is uniformly distributed over $\mathbb F_4$ except that $R_1$ is fixed to $R_1=1$. We define a palette-alternating tree code with $b$ palette sets $\Sigma_0,\ldots,\Sigma_{b-1}$ such that $\Sigma_0=\mathbb F_4$ and $\Sigma_i=\mathbb F_2$ for $i>0$. Let $x\in\mathbb F_2^{\mathbb N}$. For every $k\in\mathbb N$, define $S_k(x)=\sum_{i=1}^kR_{k+1-i}x_i$, where addition and multiplication are performed in $\mathbb F_4$ and, as usual, $\mathbb F_2$ is identified with the unique subfield of two elements in $\mathbb F_4$. The coloring function is given by
\[
T(x)_k=\begin{cases}S_k(x),&k\equiv_b0;\\\operatorname{Tr}(S_k(x)),&\text{otherwise}.\end{cases}
\]
\begin{authorclaim}{Theorem15}
The function $T$ above is a palette-alternating tree code with rate $1-\epsilon$ and distance $\delta=\Omega(\epsilon\log^{-1}(1/\epsilon))$.
\end{authorclaim}
\begin{proof}
First, observe that $T$ is indeed an online function with rate larger than $1-\epsilon$. Further Definition9 can be carried over to the more general case of palette-alternating tree codes. We turn to prove an analog to Claim10.
\begin{authorclaim}{Claim16}
Let $v$ be a depth-$n$ vertex in $T$. Let $\ell\ge1$ and $x,y\in\mathbb F_2^{\ell-1}$. Then, for every $k\in\{1,\ldots,\ell\}$ it holds that
\[
a_v(x,y,k)=\begin{cases}R_k+S_{k-1}(y-x),&n+k\equiv_b0;\\
\operatorname{Tr}(R_k+S_{k-1}(y-x)),&\text{otherwise}.\end{cases}
\]
\end{authorclaim}
\begin{proof}
Fix $k\in\{1,\ldots,\ell\}$. Assume first that $n+k\equiv_b0$. Then,
\begin{align*}
T(p_v\circ0\circ x)_{n+k}
&=\sum_{i=1}^{n+k}R_{n+k+1-i}(p_v\circ0\circ x)_i\\
&=\sum_{i=1}^nR_{n+k+1-i}(p_v)_i
+\sum_{i=n+1}^{n+k}R_{n+k+1-i}(0\circ x)_{i-n}\\
&=\sum_{i=1}^nR_{n+k+1-i}(p_v)_i
+\sum_{i=1}^kR_{k+1-i}(0\circ x)_i.
\end{align*}
Similarly
\[
T(p_v\circ1\circ y)_{n+k}
=\sum_{i=1}^nR_{n+k+1-i}(p_v)_i+\sum_{i=1}^kR_{k+1-i}(1\circ y)_i.
\]
Thus,
\begin{align*}
a_v(x,y,k)
&=\sum_{i=1}^kR_{k+1-i}(1\circ y)_i-\sum_{i=1}^kR_{k+1-i}(0\circ x)_i\\
&=R_k+\sum_{i=1}^{k-1}R_{k-i}(y-x)_i=R_k+S_{k-1}(y-x).
\end{align*}
Assume now that $n+k\not\equiv_b0$. Using that $\operatorname{Tr}$ is $\mathbb F_2$-linear,
\begin{align*}
T(p_v\circ0\circ x)_{n+k}
&=\operatorname{Tr}\left(\sum_{i=1}^{n+k}R_{n+k+1-i}(p_v\circ0\circ x)_i\right)\\
&=\operatorname{Tr}\left(\sum_{i=1}^nR_{n+k+1-i}(p_v)_i\right)
+\sum_{i=n+1}^{n+k}\operatorname{Tr}(R_{n+k+1-i})(0\circ x)_{i-n}\\
&=\operatorname{Tr}\left(\sum_{i=1}^nR_{n+k+1-i}(p_v)_i\right)
+\sum_{i=1}^k\operatorname{Tr}(R_{k+1-i})(0\circ x)_i.
\end{align*}
Similarly
\[
T(p_v\circ1\circ y)_{n+k}
=\operatorname{Tr}\left(\sum_{i=1}^nR_{n+k+1-i}(p_v)_i\right)
+\sum_{i=1}^k\operatorname{Tr}(R_{k+1-i})(1\circ y)_i.
\]
Thus, again by $\mathbb F_2$-linearity of $\operatorname{Tr}$,
\[
a_v(x,y,k)=\operatorname{Tr}(R_k)+\sum_{i=1}^{k-1}\operatorname{Tr}(R_{k-i})(y-x)_i
=\operatorname{Tr}(R_k+S_{k-1}(y-x)).
\]
\end{proof}
\begin{authorclaim}{Claim17}
Let $v$ be a depth-$n$ vertex and $x,y\in\mathbb F_2^{\ell-1}$ distinct. Then, the random variables $a_v(x,y,1),\ldots,a_v(x,y,\ell)$ are independent. Moreover, let $k\in[\ell]$. If $n+k\equiv_b0$ then $a_v(x,y,k)$ is uniformly distributed over $\mathbb F_4$ and otherwise it is uniform over $\mathbb F_2$.
\end{authorclaim}
\begin{proof}
By Claim16, if $n+k\equiv_b0$ then $a_v(x,y,k)=R_k+L_k$ where $L_k$ is a linear combination of $R_1,\ldots,R_{k-1}$. Thus, in this case, $a_v(x,y,k)$ is independent of the joint distribution of $a_v(x,y,1),\ldots,a_v(x,y,k-1)$. Otherwise, namely $n+k\not\equiv_b0$, we have that $a_v(x,y,k)=\operatorname{Tr}(R_k+L_k)=\operatorname{Tr}(R_k)+\operatorname{Tr}(L_k)$. Since for every fixing of $L_k$, $a_v(x,y,k)$ is uniform over $\mathbb F_2$, we have that $a_v(x,y,k)$ is independent of the joint distribution of $a_v(x,y,1),\ldots,a_v(x,y,k-1)$. As this holds for every $k\in[\ell]$ we have that $a_v(x,y,1),\ldots,a_v(x,y,\ell)$ are independent and their marginal distributions are as stated.
\end{proof}
\begin{authorclaim}{Claim18}
Let $u,v$ be two vertices with depth $n,m$, respectively such that $n\equiv_bm$. Let $x,y\in\mathbb F_2^{\ell-1}$. Then, $h_v(x,y)=h_u(x,y)$, and $h_v(x,y)=h_v(0^{\ell-1},y-x)$.
\end{authorclaim}
\begin{proof}
Let $C_k=R_k+S_{k-1}(y-x)$. By Claim16,
\[
a_u(x,y,k)=\begin{cases}C_k,&n+k\equiv_b0;\\\operatorname{Tr}(C_k),&\text{otherwise}.\end{cases}
\]
As $C_k$ is independent of the choice of $u$ and $n\equiv_bm$ we have that $a_u(x,y,k)$ is the same random variable as $a_v(x,y,k)$. Since this holds for every $k$, we have that $h_v(x,y)=h_u(x,y)$.
We turn to prove the the second asserted equality. Assume first that $k\in[\ell]$ is such that $n+k\equiv_b0$. By Claim16,
\begin{align*}
a_u(x,y,k)&=R_k+S_{k-1}(y-x)\\
&=R_k+S_{k-1}((y-x)-0^{\ell-1})
=a_u(0^{\ell-1},y-x,k),
\end{align*}
where observe that for the last equality we are using the fact that $\mathbb F_2$ is a subfield of $\mathbb F_4$ and so $y-x\in\mathbb F_2^{\ell-1}$. Indeed, recall that $a_v$'s second argument is a binary string and so the equality above would have been meaningless otherwise. The case $n+k\not\equiv_b0$ follows by a similar argument and using the $\mathbb F_2$-linearity of $\operatorname{Tr}$.
\end{proof}
Given Claim18, we can simplify our notation as follows. Let $v_0$ denote the root of the tree. For $i=1,\ldots,b-1$ let $v_i$ denote the left son of $v_{i-1}$. For every $i\in\{0,1,\ldots,b-1\}$ and $x\in\{0,1\}^{\ell-1}$ we define the random variables $a_i(x,k)=a_{v_i}(0^\ell,1\circ x,k)$, $h_i(x)=h_{v_i}(0^{\ell-1},x)$.
Define $\delta=c_1\epsilon\log^{-1}(1/\epsilon)$, $\ell_0=12\lceil\epsilon^{-1}\log(1/\epsilon)\rceil$, for some constant $c_1\in[0,1]$ to be set later on. Observe that for every fixing of the sequence $\{R_i\}$, $T$ is a palette-alternating tree code with distance $\delta$ if and only if for every $x\in\{0,1\}^{\ell-1}$ and $i\in\{0,1,\ldots,b-1\}$ it holds that $h_i(x)\ge\delta\ell$. Indeed, by definition, $T$ is a palette-alternating tree code with distance $\delta$ if and only if for every vertex $v$, $\ell\ge1$, and every distinct $x,y\in\{0,1\}^{\ell-1}$ it holds that $h_v(x,y)\ge\delta\ell$. However, by Claim18, the random variable $h_v(x,y)$ is the same as the random variable $h_i(y-x)$ for $i=\operatorname{depth}(v)\bmod b$.
For $x\in\{0,1\}^{\ell-1}$ and $i\in\{0,1,\ldots,b-1\}$ denote by $E_i(x)$ the event $h_i(x)<\delta\ell$. Note that as $R_1=1$ and since $\operatorname{Tr}(1)=1$ we have that $h_i(x)\ge1$ for every $x$. Thus, for $|x|<\ell_0$ we have that $h(x)/(|x|+1)\ge1/\ell_0$. Therefore, in order to prove Theorem15 it suffices to prove that $\Pr[\bigcup_{|x|\ge\ell_0}\bigcup_{i=0}^{b-1}E_i(x)]<1$.
Indeed, this will give a bound of $\min(1/\ell_0,\delta)=\Omega(\epsilon\log^{-1}(1/\epsilon))$ on the distance.
Fix $x\in\{0,1\}^{\ell-1}$ and $i\in\{0,1,\ldots,b-1\}$. Observe that $E_i(x)$ holds if and only if there exists a set $\mathcal T\subseteq[\ell]$ of size $\lceil(1-\delta)\ell\rceil$ such that for every $k\in\mathcal T$, $a_i(x,k)=0$. For ease of readability we ignore the ceiling in the calculations below. Recall that $a_i(x,1),\ldots,a_i(x,\ell)$ are independent. Further, $1-1/b$ fraction of them are uniform over $\mathbb F_2$ whereas the remaining $1/b$ fraction are uniform over $\mathbb F_4$. Note that by our choice of parameters, $\delta<1/b$. Thus, for any $\gamma\ge0$ and a fixed $\mathcal T$, we have that
\begin{align*}
\Pr[\forall k\in\mathcal T\ a_i(x,k)=0]
&\le2^{-(1-1/b-\gamma)\ell}4^{-(1/b-\delta+\gamma)\ell}\\
&\le2^{-(1-1/b)\ell}4^{-(1/b-\delta)\ell}
=2^{-(1+1/b-2\delta)\ell}.
\end{align*}
By taking the union bound over the choice of $\mathcal T$, and using Lemma7, we get that
\[
\Pr[E_i(x)]\le\binom\ell{\lceil(1-\delta)\ell\rceil}2^{-(1+1/b-2\delta)\ell}
\le2^{-(1+1/b-2\delta-H(\delta))\ell}.
\]
By the union bound,
\begin{align}
\Pr[\bigcup_{|x|\ge\ell_0}\bigcup_{i=0}^{b-1}E_i(x)]
&\le\sum_{|x|\ge\ell_0}\sum_{i=0}^{b-1}\Pr[E_i(x)]\tag{5}\\
&\le b\cdot\sum_{\ell=\ell_0}^\infty2^{\ell-1}\cdot2^{-(1+1/b-2\delta-H(\delta))\ell}\notag\\
&=\frac b2\sum_{\ell=\ell_0}^\infty2^{(H(\delta)+2\delta-1/b)\ell}.\notag
\end{align}
By taking $c_1$ sufficiently small and using Lemma8, we get that $H(\delta)+2\delta-1/b\le-\epsilon/3$. Therefore, Equation(5) is bounded above by
\[
b\sum_{\ell=\ell_0}^\infty2^{-\epsilon\ell/3}
=b\cdot\frac{2^{-\epsilon\ell_0/3}}{1-2^{-\epsilon/3}}
\le\frac{b\epsilon^4}{1-2^{-\epsilon/3}}
\le\frac{2\epsilon^3}{1-2^{-\epsilon/3}},
\]
where the penultimate inequality follows by our choice of $\ell_0$ and the last inequality follows since $b=\lceil1/\epsilon\rceil$. One can verify that the above is strictly bounded by 1 for any $\epsilon<1/3$.
\end{proof}
